\section{Introduction}

Runtime assurance is the standard way to let a controller learn without
letting it crash the plant.  A Simplex decision module, a runtime shield, or a
model-predictive safety filter holds action authority during adaptation: it
replaces unsafe policy proposals with actions from a trusted baseline, so
every transition collected while the policy trains satisfies the safety
certificate~\cite{seto_simplex_1998,alshiekh_shielding_2018,
wabersich_linear_mpsc_2018}.  Deployments do not keep that machinery forever.
Published lifecycles and ordinary engineering practice retire the protective
authority once commissioning succeeds and promote the raw policy snapshot to a
lower-latency control target, which removes permanent computation,
intervention, and dependence on a backup
service~\cite{carr_partial_shielding_2023,hsu_sim_lab_real_2023}.  A finite
release check --- a bounded-horizon prediction, a certified rollout, a
commissioning test --- authorizes that promotion.

We show that the release check does not transfer across the retirement it
authorizes.  Protection makes the executed composition safe; it does not make
the raw policy safe.  The same finite check is sound while it retains
authority and unsound once it is spent at release, because release changes the
subject being authorized, the authority enforcing the decision, and the
horizon over which the evidence must hold.  We call the underlying flaw
\emph{evidence reuse across an authority transition}: evidence collected
during protected adaptation is reused to authorize a different subject, under
a different authority, with different provenance, information, and horizon.

\begin{figure}[t]
\centering
\fbox{\begin{minipage}[c][1.74in][c]{0.92\columnwidth}
\centering
\textbf{Asynchronous learning-data plane}\\[0.25em]
reward record $\xrightarrow{\text{bounded write}}$ trusted updater
$\rightarrow$ raw candidate\\
$\hspace{3.0em}\rightarrow$ one-time check $\rightarrow$ snapshot export\\[0.35em]
\rule{0.88\columnwidth}{0.3pt}\\[-0.15em]
\textbf{High-integrity runtime-assurance plane}\\[0.25em]
authenticated state $\rightarrow$ predictor/selector $\rightarrow$ actuator\\
$\hspace{3.1em}\uparrow$ trusted baseline\\[0.35em]
resident after check: timely switch\\
authority removed: delayed failure
\end{minipage}}
\caption{Evaluated trust boundary.  The benchmark begins after bounded
reward-record access is obtained; it does not model the initial service
compromise.  Upstream provenance or recomputation removes this attack path,
whereas resident predictive authority contains its downstream physical
consequence.}
\label{fig:lifecycle-placeholder}
\end{figure}

The gap is measurable without an adversary.  In our controlled study,
snapshots that passed a five-step release check failed on 2 of 120
snapshot--state pairs once released raw, whereas the identical prediction
retained at runtime failed on none.  Residence alone is not the remedy: on a
disjoint cohort, a one-step resident kernel failed on 14 of 71 accepted pairs
and reached an empty kernel 41 times.  What separates these outcomes is
whether a check with enough lookahead to leave switching lead time keeps
authority at runtime.

An upstream adversary widens the gap rather than creating it.  While the
shield masks unsafe proposals, corrupted update data steers the learner
without producing a single physical violation, and the damage stays latent
until authority is removed: bounded reward-record edits raised the same
120-pair count from 2 to 27.  We call this \emph{shield-retirement poisoning}.
Its applicability is narrow by construction (Section~\ref{sec:threat}): where
the runtime plane can recompute reward from authenticated transitions, trusted
recomputation removes the channel outright, as our frozen integrity audit
confirms.  The contract gap survives that defense.

The attacker we evaluate is deliberately weak.  Operational technology (OT)
guidance separates runtime-control assets from the data-management services
and historians that exchange information across
zones~\cite{nist_ot_security_2023,mitre_data_historian}, and our threat model
begins after compromise of a reward-ingestion principal in that learning-data
plane.  This is an explicit split-trust precondition: the attacker may modify
bounded scalar reward records consumed by a policy-gradient updater, and
nothing else.  Authenticated runtime states, actions, learner code,
parameters, optimizer state, certificates, and deployment states remain
outside its write authority (Fig.~\ref{fig:lifecycle-placeholder}).

Existing mechanisms address neighboring problems.  Simplex and its neural and
barrier-based successors keep a decision module resident
\cite{seto_simplex_1998,phan_neural_simplex_2020,damare_bb_simplex_2025}.
Policy repair and safe projection constrain a learned policy before deployment
\cite{zhou_runtime_policy_repair_2020,lu_repair_preservation_2024,
chow_safe_policy_2021}.  Reward-poisoning work studies how adversarial
learning signals redirect updates, but not behind a shield.  None treats the
release decision itself as the object of study: the unresolved composition is
a transition in which evidence produced under retained authority is spent
once, at release, to authorize a raw policy that will run without it.

\subsection{Evidence and Contributions}

We separate the contract gap from the channel that widens it.  The controlled
Safe-Control-Gym study supplies the counts above under a five-step check that
accepted every candidate, and adds a locked horizon sweep placing the usable
lookahead between an insufficient one-step check and an over-conservative
twenty-step one.  The channel is not confined to our harness: Carr et al.'s
published workflow trains a policy behind a Storm belief-support shield and
later retires the shield's authority~\cite{carr_partial_shielding_2023}, and
adding bounded reward-record poisoning only during the protected bootstrap
increased raw-policy violations at retirement by factors of 3.6 and 1.6 in two
of three provenance-complete paired REINFORCE training runs.

Our contributions are:

\begin{itemize}
    \item We formalize release-time acceptance as an evidence tuple over
    subject, authority, provenance, information, and horizon, and show that a
    deployment must re-establish every component the authority transition
    changes.
    \item We measure the resulting gap without an adversary, separating
    one-time raw release from resident predictive authority, one-step
    filtering, and commit admission, and locating the lookahead at which a
    retained check is neither insufficient nor over-conservative.
    \item We instantiate one upstream channel that widens the gap ---
    bounded reward-record writes --- in the controlled study and in a
    third-party shield-retirement lifecycle, and report the integrity controls
    that close that channel while leaving the gap open.
    \item We characterize the attack's empirical boundary across systems,
    learners, integrity controls, and deployment contracts: high
    retirement-state disagreement excludes every observed positive outcome in
    the locked corpus, while low disagreement defines an opportunity region
    without guaranteeing success.
\end{itemize}
